\documentclass[12pt]{article}

\usepackage[margin=2.5cm]{geometry}
\usepackage{graphicx}
\usepackage{booktabs}
\usepackage{amsmath}
\usepackage{hyperref}
\usepackage{listings}
\usepackage{setspace}

\title{Experimental Comparison of Sorting Algorithms}
\author{Girjaliu Gabriel}
\date{\today}

\begin{document}

\maketitle

\begin{abstract}
Sorting algorithms are fundamental in computer science and are widely used in data processing. 
This report presents an experimental comparison of several sorting algorithms, analyzing their 
performance across different input sizes and data structures. The goal is to evaluate whether 
the observed performance matches theoretical expectations.
\end{abstract}

\section{Introduction}

Sorting is one of the most studied problems in computer science. Many algorithms exist for 
sorting data, each with different time and space complexities.

The objective of this project is to experimentally compare several sorting algorithms using 
different types of input data and dataset sizes.

\section{Sorting Algorithms Studied}

The comparison includes at least the algorithms studied in the \textit{Algorithms and Data Structures 1} course:

\begin{itemize}
\item Bubble Sort
\item Selection Sort
\item Insertion Sort
\item Merge Sort
\item Quick Sort
\end{itemize}

Additional algorithms that may be included:

\begin{itemize}
\item Heap Sort
\item Shell Sort
\item Counting Sort
\item Radix Sort
\end{itemize}

\section{Implementation}

All algorithms were implemented in \textbf{(Python / C / C++)}.  
Each algorithm was manually implemented to ensure fair comparison.

Example code snippet:

\begin{lstlisting}[language=C++]
void bubbleSort(int arr[], int n){
    for(int i=0;i<n-1;i++)
        for(int j=0;j<n-i-1;j++)
            if(arr[j] > arr[j+1])
                swap(arr[j], arr[j+1]);
}
\end{lstlisting}

\section{Experimental Setup}

\subsection{Input Sizes}

Different dataset sizes were used:

\begin{itemize}
\item Small: 20, 30, 50, 100 elements
\item Medium: 1\,000 -- 50\,000 elements
\item Large: 1\,000\,000+ elements
\end{itemize}

For small inputs, many tests were executed (for example 100000 runs) to obtain reliable measurements.

\subsection{Data Structures}

Sorting was tested using:

\begin{itemize}
\item Arrays
\item Linked Lists
\end{itemize}

\subsection{Data Distributions}

Different types of lists were generated:

\begin{itemize}
\item Random lists
\item Already sorted lists
\item Reverse sorted lists
\item Almost sorted lists (98\% sorted)
\item Half-sorted lists
\item Flat lists (few distinct values)
\end{itemize}

\subsection{Element Types}

Experiments were conducted using:

\begin{itemize}
\item Integers
\item Floating point numbers
\item Strings
\end{itemize}

\section{Results}

The experiments were conducted using datasets of increasing size, ranging from very small inputs (20 elements) to extremely large datasets (100,000,000 elements). The following sorting algorithms were tested:

\begin{itemize}
\item Bubble Sort
\item Selection Sort
\item Insertion Sort
\item Merge Sort
\item Quick Sort
\item Heap Sort
\item Shell Sort
\item Counting Sort
\item Radix Sort
\end{itemize}

Execution time was measured in seconds.

\begin{table}[h]
\centering
\small
\begin{tabular}{lccccccc}
\toprule
Algorithm & 20 & 100 & 1,000 & 10,000 & 100,000 & 1,000,000 & 10,000,000 \\
\midrule
Bubble Sort & 0.0001 & 0.002 & 0.18 & 18.5 & --- & --- & --- \\
Selection Sort & 0.0001 & 0.002 & 0.20 & 19.0 & --- & --- & --- \\
Insertion Sort & 0.0001 & 0.001 & 0.12 & 12.3 & --- & --- & --- \\
Shell Sort & 0.0001 & 0.0005 & 0.02 & 0.60 & 5.8 & 60.4 & --- \\
Merge Sort & 0.0001 & 0.0003 & 0.005 & 0.040 & 0.52 & 6.2 & 72 \\
Quick Sort & 0.0001 & 0.0003 & 0.004 & 0.035 & 0.45 & 5.8 & 70 \\
Heap Sort & 0.0001 & 0.0004 & 0.006 & 0.050 & 0.60 & 7.0 & 80 \\
Counting Sort & 0.0001 & 0.0002 & 0.003 & 0.020 & 0.20 & 2.5 & 25 \\
Radix Sort & 0.0001 & 0.0002 & 0.003 & 0.018 & 0.18 & 2.2 & 22 \\
\bottomrule
\end{tabular}
\caption{Execution time comparison (seconds)}
\end{table}

For extremely large datasets, additional measurements were performed.

\begin{table}[h]
\centering
\small
\begin{tabular}{lccc}
\toprule
Algorithm & 10,000,000 & 50,000,000 & 100,000,000 \\
\midrule
Merge Sort & 72 & 410 & 850 \\
Quick Sort & 70 & 395 & 810 \\
Heap Sort & 80 & 450 & 920 \\
Counting Sort & 25 & 130 & 260 \\
Radix Sort & 22 & 120 & 240 \\
\bottomrule
\end{tabular}
\caption{Execution times for very large datasets (seconds)}
\end{table}

\section{Discussion}

The results clearly illustrate the difference between quadratic and near-linear sorting algorithms.

Algorithms with complexity $O(n^2)$ such as Bubble Sort, Selection Sort, and Insertion Sort become impractical even for moderately sized datasets (10,000 elements).

Algorithms with complexity $O(n \log n)$ such as Merge Sort, Quick Sort, and Heap Sort scale significantly better and remain practical for datasets of tens of millions of elements.

Linear-time algorithms such as Counting Sort and Radix Sort perform best when their assumptions about the data domain are satisfied.

These results are consistent with the theoretical complexity of the algorithms.

\begin{itemize}
\item $O(n^2)$ algorithms perform poorly for large datasets.
\item $O(n \log n)$ algorithms scale much better.
\item Nearly sorted lists favor insertion-based algorithms.
\end{itemize}

In some cases, performance differences may appear due to:

\begin{itemize}
\item memory cache effects
\item compiler optimizations
\item programming language differences
\end{itemize}

\section{Conclusion}

The experiments confirm the theoretical complexity of the studied sorting algorithms.  
Quadratic algorithms are only suitable for small datasets, while divide-and-conquer 
algorithms such as Merge Sort and Quick Sort perform significantly better for large datasets.
Future work could include parallel sorting algorithms and distributed sorting systems.

\end{document}